\section{Architecture}

The implementation is a Rust workspace with seven crates whose dependency order is
\texttt{effects} $\rightarrow$ \texttt{guards} $\rightarrow$ \texttt{cache} $\rightarrow$
\texttt{trace} $\rightarrow$ \texttt{sandbox} $\rightarrow$ \texttt{core} $\rightarrow$
\texttt{cli}. The order is a correctness device: nothing that can decide to reuse is reachable from
anything that only records.

\paragraph{Tracing.}
A command is launched under \texttt{ptrace} and every modelled syscall is decoded into an effect
with its arguments, result and errno. Children and threads are followed. The trace is bounded by
the process tree, not by wall-clock time.

\paragraph{Effect classification.}
The classifier folds the effect list into a determinism class. The important property is that
classification is a function of the \emph{observed} effects only: a syscall the tracer does not
model cannot be silently dropped, because an unmodelled syscall is itself an effect that forces
refusal. This is what makes the refusal list auditable, and it is also what makes some real
programs permanently ineligible (Section~\ref{sec:negative}).

\paragraph{Guard construction.}
Each read becomes a guard that re-opens the recorded pathname, re-reads it and compares a BLAKE3
digest~\cite{blake3}; each metadata observation becomes a guard over the recorded fields; each
access check becomes a guard over the recorded answer, including a negative answer. Guards replay
the pathname the process used: there is no lexical \texttt{..} collapse and no extra symlink
resolution, because either would change which file the process would see on the next run.

\paragraph{Identity and candidate lookup.}
The lookup key is the execution identity of Section~\ref{sec:model}. A bounded number of alternate
candidates is retained per identity. A candidate whose contract or class does not match is skipped
without a trace penalty.

\paragraph{Cache and CAS.}
Objects are content-addressed with BLAKE3~\cite{blake3} and stored on disk with SQLite for the
record index. Before any cached byte is exposed, the object is re-hashed; a mismatch is a
\emph{corruption} result that discards the candidate and retraces. Corruption therefore costs a
re-execution, never a wrong answer.

\paragraph{Validated replay.}
On a hit the sequence is: load stream objects and verify them, run the single guard pass, stage and
verify output objects, restore outputs, reproduce any inherited stream offsets, then expose the
streams. The type system enforces the order: the cached bytes are held in an unforgeable value
that only a passing guard pass can convert into an exposable one. This is a structural guarantee,
not a convention.

\paragraph{Transactional restoration.}
Outputs are staged and verified before any destination is touched. A single output is renamed into
place. Multiple outputs carry a durable restore intent so that a soft failure rolls back and a hard
kill causes the next attempt to refuse a hit rather than expose a partially restored set.

\paragraph{Deoptimisation.}
A failing guard, a corrupt object, or a failed output precondition discards the candidate and
re-executes, storing the new execution. Deoptimisation is a first-class outcome with its own
report line, because a cache that cannot say ``I was wrong'' cannot be audited.

\paragraph{Enforcement.}
With \texttt{--enforce}, a re-executed record receives a Landlock ruleset built from its recorded
dependencies and a seccomp filter that denies network syscalls and a small set of privileged ones.
The child verifies the filter in itself (\texttt{PR\_GET\_SECCOMP}, \texttt{PR\_GET\_NO\_NEW\_PRIVS})
before the command runs. Only then is the record labelled \proven{}. The label records
installation, not equivalence.
